\uuid{PoCf}
\titre{Équation droite/plan}
\auteur{bodin ; hanani ; mzari}
\organisation{exo7}
\licence{CC BY-NC-SA 4.0}
\datecreate{2020-01-20}
\niveau{L1}
\module{Géométrie}
\chapitre{Géométrie de l'espace}
\sousChapitre{Plans}
\difficulte{2}
\theme{Plans -- Droites, équation droite, plan}
\exo7qcmid{lille-1.geometrie-espace.013}

\contenu{
\texte{
Pour ces questions, l'espace est muni d'un repère orthonormé direct  $(O, \vec{i}, \vec{j}, \vec{k})$.
}

\begin{question}
\qclassification{141.03,141.04}
\qtags{motcle=équation droite/plan}
On considère la droite $D : \left\{\begin{array}{ccl}x&=&1-t\\y&=&t\\ z&=&-1+2t,\quad (t \in \Rr) \end{array}\right.$ et le plan $P$  passant par $A(0,1,1)$ et perpendiculaire à $D$.  Quelles sont les assertions vraies ?
\begin{answers}

    \good{Une équation cartésienne de $P$ est $x-y-2z+3=0$.}

    \bad{Une équation cartésienne de $P$ est $x-2y-2z+2=0$.}

   \good{Une représentation paramétrique de $P$ est :
 $$\left\{\begin{array}{ccl}x&=&t\\y&=&3+t-2s\\ z&=&s,\quad (t,s \in \Rr) \end{array}\right.$$}

    \bad{Une représentation paramétrique de $P$ est :
 $$\left\{\begin{array}{ccl}x&=&t\\y&=&2+t-2s\\ z&=&s,\quad (t,s \in \Rr) \end{array}\right.$$}

\end{answers}
\begin{explanations}
Un vecteur directeur de $D$ est un vecteur normal à $P$.
\end{explanations}
\end{question}
}
